% ==============================================================================
% CAPÍTULO 11: ASPECTOS DE PUESTA EN MARCHA DEL SOFTWARE (Orquestador del Capítulo)
% ==============================================================================
\preportadacapitulo{11}{Aspectos de Puesta en Marcha del Software}

\newpage
\section{Capítulo 11: Aspectos de Puesta en Marcha del Software}\label{cap:puesta_marcha}

% [INSTRUCCIÓN]: Puesta en marcha y despliegue (fase de Transición del PUDS). Contenido propio de esta versión.
% [VERIFICAR]: las normas citadas se leyeron en LexiVox el 19-nov-2026 y se contrastan con la Gaceta Oficial antes de la entrega; los precios de servidores y de la tienda de aplicaciones se toman de las páginas oficiales en la fecha de cotización.

Este capítulo corresponde a la fase de Transición del Proceso Unificado: lleva el sistema a producción y define las condiciones bajo las que se entrega. Se comparan las plataformas de alojamiento, se definen las licencias, y se resumen los términos, condiciones y políticas de privacidad.

\subsection{11.1. Aspectos de Puesta en Marcha y Despliegue}

El sistema tiene cuatro piezas que se despliegan: la API, el módulo de inteligencia artificial, la base de datos con el almacenamiento de archivos, y las interfaces (panel web y aplicación móvil Android). Las primeras se empaquetan con Docker, como en el entorno de desarrollo (Capítulo~\ref{cap:infraestructura}).

\subsubsection{11.1.1. Amazon Web Services (AWS)}

AWS es la plataforma de nube elegida para el sistema (sección 1.11). Permite alojar los contenedores de la API y del módulo de IA, la base de datos PostgreSQL como servicio administrado y el almacenamiento de archivos. Ofrece copias de seguridad, escalamiento y monitoreo. Su uso permite aumentar los recursos cuando crece la cantidad de postulaciones durante la convocatoria.

\subsubsection{11.1.2. Microsoft Azure}

Azure ofrece servicios equivalentes (máquinas virtuales, contenedores y bases de datos administradas) y se integra bien con entornos que ya usan productos de Microsoft. Se la considera una alternativa si la institución cliente tuviera ya un contrato con ella.

\subsubsection{11.1.3. Google Cloud Platform (GCP)}

GCP ofrece también contenedores, bases de datos administradas y almacenamiento, con buenas herramientas de análisis de datos. Es la tercera alternativa comparada.

\subsubsection{11.1.4. Play Store (despliegue móvil)}

La aplicación móvil del postulante tiene como prioridad la plataforma Android (sección 1.11). Se publica en Google Play Store con una cuenta de desarrollador, un paquete firmado, una ficha con descripción y capturas, y la política de privacidad de la sección 11.5. Google revisa la aplicación antes de publicarla. Las actualizaciones posteriores se suben como nuevas versiones.

\subsection{11.2. Planes y Cotización (Resumen de Servidores)}

La tabla compara las tres plataformas en los servicios que necesita el sistema.

\begin{table}[H]
\centering
\caption{Servicios que requiere el sistema en cada plataforma de nube}
\label{tab:nube}
\small
\begin{tabular}{@{}p{3.6cm}p{3.6cm}p{3.6cm}p{3.6cm}@{}}
\toprule
\textbf{Necesidad} & \textbf{AWS} & \textbf{Azure} & \textbf{GCP} \\
\midrule
Contenedores de la API y de la IA & Servicio de contenedores o máquina virtual & Servicio de contenedores o máquina virtual & Servicio de contenedores o máquina virtual \\
Base de datos PostgreSQL & Servicio administrado & Servicio administrado & Servicio administrado \\
Archivos (documentos y actas) & Almacenamiento de objetos & Almacenamiento de objetos & Almacenamiento de objetos \\
Costo mensual estimado & Unos USD 100 (máquina virtual y disco; ver el texto siguiente) & No se cotizó & No se cotizó \\
\bottomrule
\end{tabular}

\vspace{2pt}
\footnotesize\textit{Nota.} Elaboración propia. El costo de AWS se calculó con los precios de la lista oficial de AWS (API de precios) el 19 de noviembre de 2026; Azure y GCP no se cotizaron porque el sistema se desplegó en AWS.
\end{table}

\paragraph{11.2.1. Costo del despliegue actual en AWS.}

El despliegue actual del sistema, hecho para la demostración, corre en una sola máquina virtual de Amazon EC2 en la región de São Paulo (\texttt{sa-east-1}), de tipo \texttt{c7i-flex.large} (2 vCPU y 4 GiB de memoria), con un disco de 30 GB de tipo gp3. Dentro de la máquina, el sistema corre con Docker: el servidor web, el panel, el sitio de NEXWARE, la API, el módulo de IA, PostgreSQL y Redis. Las imágenes se guardan en Amazon ECR (cuatro repositorios), los archivos de los postulantes en Amazon S3, el correo sale por Amazon SES y el acceso público, con HTTPS, pasa por Amazon CloudFront. Es decir, la base de datos no es un servicio administrado, como se plantea para una producción definitiva.

Con los precios de la lista oficial de AWS para esa región, la máquina virtual cuesta USD 0,13067 por hora, es decir, unos USD 95,39 al mes si está encendida todo el mes (730 horas), y el disco cuesta USD 0,152 por GB al mes, es decir, USD 4,56. Juntos suman unos USD 100 mensuales. Esta cifra no incluye el almacenamiento y las transferencias de S3, CloudFront, ECR, SES ni el cobro por la dirección IP pública, que dependen del uso real y no se calcularon. La cuenta no tiene activado el Explorador de Costos, por lo que no hay una factura real que confirmar. Esta estimación corresponde a la demostración y no a la carga prevista de producción (20,000 postulaciones por gestión y al menos 50 administrativos concurrentes).

% [PENDIENTE]: costo real de AWS cuando se active el Explorador de Costos. Azure y GCP no se cotizan (decisión del equipo: el sistema usa AWS).

\subsection{11.3. Tipos de Licencia Ofrecidos}

NEXWARE ofrece el sistema con una licencia de uso, no con la venta del código. Se consideran dos modalidades:

\begin{itemize}
    \item \textbf{Licencia de uso con vigencia anual:} la institución usa el sistema durante un año y recibe las actualizaciones de mantenimiento. Se renueva cada año.
    \item \textbf{Licencia de uso permanente:} la institución usa la versión entregada sin límite de tiempo; las actualizaciones y el soporte se contratan aparte.
\end{itemize}

En ambos casos el código fuente sigue siendo de NEXWARE (sección 5.2). La institución no puede ceder ni revender el sistema. Los precios se definen en la sección 8.7.3.

\subsection{11.4. Términos y Condiciones}

\subsubsection{11.4.1. Normas bolivianas aplicables}

Los términos y condiciones del servicio deben ajustarse a estas normas:

\begin{itemize}
    \item \textbf{Código de Comercio} \parencite{codigocomercio1977}: regula los contratos mercantiles.
    \item \textbf{Ley N.º~1322 de Derecho de Autor} \parencite{ley1322}: protege el software y fija los derechos del titular.
    \item \textbf{Ley N.º~164 General de Telecomunicaciones, Tecnologías de Información y Comunicación (2011):} marco general de los servicios digitales, el comercio electrónico y la firma digital.
    \item \textbf{Ley N.º~453 General de los Derechos de las Usuarias y los Usuarios y de las Consumidoras y los Consumidores (2013):} derechos de los usuarios frente a quien presta un servicio.
\end{itemize}

\subsubsection{11.4.2. Requisitos para plataformas con suscripciones digitales en Bolivia}

Para ofrecer el sistema con licencia de pago periódico, NEXWARE debe:

\begin{itemize}
    \item Estar inscrita en el Registro de Comercio y en el padrón del SIN, y emitir factura por cada pago (sección 5.1.4).
    \item Dar a conocer al cliente, antes de contratar, el precio, la vigencia, la forma de renovación y de cancelación.
    \item Atender los reclamos del cliente y mantener un canal de soporte (Capítulo~\ref{cap:sitio_web}).
    \item Firmar un contrato con el alcance del servicio, las responsabilidades de cada parte y las condiciones de la licencia.
\end{itemize}

\subsection{11.5. Políticas de Seguridad y Privacidad}

El sistema trata datos personales y socioeconómicos de estudiantes (ingresos familiares, dependientes, patrimonio y documentos). Esa información exige especial cuidado.

\subsubsection{11.5.1. Leyes y normas bolivianas aplicables sobre datos personales}

\begin{itemize}
    \item \textbf{Constitución Política del Estado:} reconoce el derecho a la privacidad, la intimidad, la honra, el honor, la propia imagen y la dignidad (artículo 21, numeral 2) y establece la acción de protección de privacidad para quien vea afectado su derecho a la intimidad y la privacidad por datos registrados en archivos o bancos de datos públicos o privados (artículo 130).
    \item \textbf{Ley N.º~164 (2011):} reconoce a los usuarios de telecomunicaciones el derecho a exigir respeto a la privacidad e inviolabilidad de sus comunicaciones y a recibir del proveedor protección de sus datos personales contra la publicidad no autorizada (artículo 54, numerales 6 y 17); el artículo 56 obliga a garantizar la inviolabilidad y el secreto de las comunicaciones y la protección de los datos personales.
    \item \textbf{Ley N.º~1080 de Ciudadanía Digital (2018):} establece las condiciones y responsabilidades para el acceso pleno y el ejercicio de la ciudadanía digital (artículo 1), lo que incluye la identidad y los trámites digitales.
\end{itemize}

Hasta donde el equipo ha podido verificar, Bolivia no cuenta con una ley general de protección de datos personales: la Agencia de Gobierno Electrónico y Tecnologías de Información y Comunicación (AGETIC) presentó un anteproyecto en 2024 que, según las fuentes consultadas, seguía sin promulgarse en 2026. Por eso el sistema aplica, además de las normas anteriores, buenas prácticas reconocidas.

% [VERIFICAR]: los artículos citados se leyeron el 19-nov-2026 en LexiVox (no es la fuente oficial); contrastarlos con la Gaceta Oficial de Bolivia antes de la entrega. El estado del anteproyecto de protección de datos viene de fuentes secundarias (AGETIC y organizaciones de la sociedad civil); confirmarlo en la Asamblea Legislativa.

\subsubsection{11.5.2. Redacción de las políticas de privacidad y gestión de datos}

Las políticas del sistema se resumen así:

\begin{itemize}
    \item \textbf{Datos recogidos:} solo los necesarios para evaluar la postulación y controlar la asistencia.
    \item \textbf{Uso:} calcular el puntaje, auditar y adjudicar. No se venden ni se ceden a terceros.
    \item \textbf{Acceso:} cada rol ve solo lo que necesita; las acciones críticas quedan en la bitácora de auditoría.
    \item \textbf{Protección:} comunicación cifrada (HTTPS), contraseñas guardadas como \textit{hash}, sesiones con tokens que vencen (Capítulo~\ref{cap:pruebas}).
    \item \textbf{Conservación:} los datos se conservan mientras dure la gestión y los plazos que fije la institución.
    \item \textbf{Derechos del titular:} el estudiante puede consultar y pedir corregir sus datos.
\end{itemize}

El texto completo de la política se publica en el sitio de NEXWARE y en la ficha de la aplicación en Play Store.

\subsection{11.6. Anexos}

\subsubsection{11.6.1. Herramientas de proyección de infraestructura (Cloud y Móvil)}

Para calcular el costo de la nube en AWS se usa la API de precios y la calculadora oficial (AWS Pricing Calculator). En el caso del despliegue móvil, la cuenta de desarrollador de Google Play tiene un pago de registro que se consulta en la página oficial de Google.

\subsubsection{11.6.2. Marco legal y normativo (Bolivia)}

Las normas del capítulo son el Código de Comercio, la Ley N.º~1322, la Ley N.º~164, la Ley N.º~453, la Ley N.º~1080 y la Constitución Política del Estado. Los trámites de constitución se tratan en el Capítulo~\ref{cap:legal}.

\subsubsection{11.6.3. Recursos y evidencia del proyecto}

El sistema está desplegado y se accede a él en las direcciones siguientes:

\begin{itemize}
    \item \textbf{Panel web del sistema:} \url{https://d38nteg7p7v4qz.cloudfront.net/login}
    \item \textbf{Sitio web de NEXWARE:} \url{https://d38nteg7p7v4qz.cloudfront.net/landing-page/es/}
\end{itemize}

El despliegue usa los servicios de AWS descritos en la sección 11.2.1. Para el despliegue de la aplicación móvil, el equipo prepara un archivo de instalación (APK) de Android; todavía no está publicada en Google Play Store.

% [PENDIENTE]: enlace al APK (Drive), captura del despliegue y, si se decide, enlace público del repositorio. No incluir evidencias que no existan.
